\section{Відповіді на контрольні питання}

У цьому розділі наведено теоретичні відповіді на контрольні питання з метою 
підтвердження засвоєння фундаментальних концепцій об'єктно-орієнтованого програмування.

\begin{description}
    \item[1. Що таке програма з погляду об'єктно-орієнтованої парадигми?] ~\\
    З погляду об'єктно-орієнтованої парадигми (ООП), програма є \textbf{сукупністю взаємодіючих об'єктів}, 
    кожен з яких є екземпляром певного класу. Замість послідовного виконання команд (як у процедурному 
    підході), виконання програми базується на надсиланні повідомлень між об'єктами та зміні їхніх 
    внутрішніх станів.

    \item[2. Що таке об'єкт?] ~\\
    Об'єкт — це конкретний будівельний блок об'єктно-орієнтованої програми, який 
    є \textbf{реальним екземпляром абстрактного класу} та поєднує в собі фіксовані дані 
    (стан) і методи для їхньої обробки (поведінку).

    \item[3. З яких частин складається об'єкт?] ~\\
    Кожен об'єкт в ООП складається з трьох обов'язкових частин:
    \begin{itemize}
        \item \textbf{Ім'я (ідентифікатор):} унікальна назва або посилання на об'єкт у пам'яті програми (наприклад, посилання \texttt{stackCart}).
        \item \textbf{Стан (дані/атрибути):} поточні значення змінних екземпляра, які характеризують властивості об'єкта в даний момент часу.
        \item \textbf{Поведінка (методи):} набір операцій, які об'єкт може виконувати або які можна виконати над ним.
    \end{itemize}

    \item[4. Що таке клас?] ~\\
    Клас — це \textbf{абстрактний шаблон або креслення}, на основі якого створюються реальні 
    об'єкти (екземпляри). Клас визначає новий тип даних, описуючи структуру полів та алгоритми
     методів, що будуть спільними для всіх його об'єктів.

    \item[5. З яких частин складається клас?] ~\\
    Тіло класу складається з двох основних частин, які узагальнено називаються 
    \textbf{членами класу}:
    \begin{itemize}
        \item \textbf{Поля класу (змінні екземпляра):} визначають структуру даних об'єкта.
        \item \textbf{Методи класу (функції):} містять програмний код для обробки полів та визначають зовнішній інтерфейс взаємодії.
    \end{itemize}

    \item[6. Як створити об'єкт якогось класу?] ~\\
    У мові Java об'єкт створюється динамічно у кучі (\textit{Heap}) за допомогою 
    \textbf{оператора \texttt{new}}, після якого викликається конструктор відповідного класу
     (наприклад, \texttt{Cart cart = new StackCart(8);}). Оператор \texttt{new} виділяє пам'ять 
     під об'єкт і повертає посилання на нього.

    \item[7. Що таке конструктор? Що таке конструктор за замовчуванням?] ~\\
    \textbf{Конструктор} — це спеціальний метод класу, який автоматично викликається під час створення об'єкта для ініціалізації його початкового стану. Його ім'я збігається з ім'ям класу, і він не має типу повертаємого значення [1].\\
    \textbf{Конструктор за замовчуванням} — це конструктор без параметрів, який автоматично створюється компілятором Java, якщо в класі не оголошено жодного кастомного конструктора. Він ініціалізує числові поля нулями, логічні — значенням \texttt{false}, а посилання — значенням \texttt{null}.

    \item[8. Що таке інкапсуляція? Навіщо вона потрібна? Як реалізовано інкапсуляцію в мові Java?] ~\\
    \textbf{Інкапсуляція} — це принцип ООП, що полягає в об'єднанні даних та методів їхньої обробки всередині об'єкта, а також у \textbf{прихованні внутрішньої реалізації} від прямого зовнішнього втручання [1].\\
    \textbf{Навіщо потрібна:} вона захищає дані від некоректної модифікації, ізолює контракт взаємодії від внутрішнього устрою, що дозволяє змінювати код всередині класу (наприклад, замінити стек чергою) без збоїв у решті програми [1].\\
    \textbf{Реалізація в Java:} забезпечується за допомогою \textbf{модифікаторів доступу} (\texttt{private}, \texttt{public}, \texttt{protected}). Поля класу оголошуються як \texttt{private}, а доступ до них та валідація контролюються через відкриті (\texttt{public}) методи — геттери, сеттери та кастомні методи бізнес-логіки.
\end{description}
